\subsection{Reflection and the two classical hierarchies}\label{sec:reflection}

The operator families met most often in applications are smaller than
the full monic family, and \(N_p\) and \(M_p\) do not apply to them
directly. An evolution equation that contains only odd-order
derivatives, such as the KdV, Kawahara and seventh-order KdV
equations, gives an even polynomial \(P\) of even degree. An equation
whose linear part contains only even-order derivatives, such as the
Kuramoto--Sivashinsky equation without dispersion, gives a polynomial
of odd degree whose nonconstant part is odd. Both families are the
fixed locus of a reflection, and the count inside them is again a
binomial coefficient.

We call a monic polynomial \(P\) of degree \(p\)
\emph{reflection-symmetric} if \(P-P(0)\) has the parity of
\(\rho^p\):
\begin{equation}\label{eq:symmetric-class}
 P(-\rho)=P(\rho)\quad(p\text{ even}),\qquad
 P(\rho)+P(-\rho)=2P(0)\quad(p\text{ odd}).
\end{equation}
In the evolution equation of Section~\ref{sec:intro} this means
\(b_j=0\) whenever \(p-j\) is odd. The equation is then invariant under
\((x,t)\mapsto(-x,-t)\) for even \(p\), and under
\((x,u)\mapsto(-x,-u)\) for odd \(p\).

For a monic polynomial \(A\) of degree \(d=p-1\) put
\begin{equation}\label{eq:iota}
 (\iota A)(\rho)=(-1)^dA(-\rho),\qquad\text{that is,}\qquad
 \alpha_j\longmapsto(-1)^j\alpha_j .
\end{equation}

\begin{proposition}[Reflection]\label{prop:reflection}
For every monic \(A\) of degree \(d\),
\begin{equation}\label{eq:reflection-S}
 S_{\iota A}(\rho)=-S_A(-\rho)-2A(0)A(-\rho),\qquad
 R_{\iota A}(\rho)=-R_A(-\rho).
\end{equation}
Consequently \(\iota\) is an involution of the scheme \(\Xp\).
At a point of \(\Xp\),
\begin{equation}\label{eq:reflection-pair}
 v_{\iota A}(z)=(-1)^p\bigl(v_A(-z)-s_pA(0)\bigr),\qquad
 P_{\iota A}(\rho)=(-1)^p\bigl(P_A(-\rho)+s_pA(0)\bigr).
\end{equation}
\end{proposition}

\begin{proof}
Since \(Q(-z)=1-Q(z)\), and \(D\) changes sign under \(z\mapsto-z\),
\[
 \bigl(A(D)Q\bigr)(-z)=A(0)-A(-D)Q=A(0)-(-1)^d(\iota A)(D)Q .
\]
This is the first formula in \eqref{eq:reflection-pair}.
Evaluate \eqref{eq:discrete-square} at \(-z\), and use
\eqref{eq:discrete-square} for \(\iota A\) on the left:
\[
 A(0)^2-2A(0)A(-D)Q-S_{\iota A}(D)Q=-S_A(0)+S_A(-D)Q .
\]
Every \(B(D)Q\) is a polynomial in \(Q\) without constant term, so
letting \(z\to-\infty\) gives \(S_A(0)=-A(0)^2\). The remaining terms,
and the injectivity of \(B\mapsto B(D)Q\), give the first identity in
\eqref{eq:reflection-S}. Now substitute \(S_A=AC_A+R_A\) at \(-\rho\):
\[
 S_{\iota A}(\rho)
 =(\iota A)(\rho)\,(-1)^p\bigl(C_A(-\rho)+2A(0)\bigr)-R_A(-\rho).
\]
Division by the monic polynomial \(\iota A\) is unique. Hence
\[
 R_{\iota A}(\rho)=-R_A(-\rho),\qquad
 C_{\iota A}(\rho)=(-1)^p\bigl(C_A(-\rho)+2A(0)\bigr).
\]
The first relation reads \(R_r(\iota\alpha)=(-1)^{r+1}R_r(\alpha)\), so
the linear map \eqref{eq:iota} preserves the ideal of the matching
equations. The second relation and \eqref{eq:operator-A} give the
formula for \(P_{\iota A}\).
\end{proof}

Thus \(\iota\) reads a wave in the opposite direction and moves its
zero endpoint back to the left; for odd \(p\) it also changes the sign
of the profile. The pairs \(A\) and \(\iota A\) are counted separately
in \(N_p\). For \(p=3\) the involution is \((a,b)\mapsto(-a,b)\) in
Table~\ref{tab:ks}.

Let
\[
 \Xp^{\mathrm{sym}}=\Xp\cap\{\alpha_j=0\text{ for odd }j\}
\]
be the fixed locus of \(\iota\), with its induced scheme structure, and
put
\begin{equation}\label{eq:Fp}
 F_p=\binom{p-1}{\lfloor (p-1)/2\rfloor}.
\end{equation}

\begin{theorem}[Count in the symmetric families]\label{thm:symmetric}
\begin{enumerate}
\item[\textup{(i)}] A point \(A\) of \(\Xp\) lies in
\(\Xp^{\mathrm{sym}}\) if and only if the operator \(P_A\) is
reflection-symmetric. Thus \(\Xp^{\mathrm{sym}}\) is the set of
normalized rational-exponential pairs of the family
\eqref{eq:symmetric-class}.
\item[\textup{(ii)}] If \((\mathrm R_{p-1})\) holds, then
\(\len\Xp^{\mathrm{sym}}=F_p\).
\item[\textup{(iii)}] If, in addition, \(\Xp\) is reduced, then
\(\Xp^{\mathrm{sym}}\) consists of \(F_p\) distinct pairs, each a simple
solution of the matching equations of the symmetric family. For even
\(p\) all of them are pulses, and for odd \(p\) all of them are fronts.
If \(2p-1\) is prime, they are the points whose label in
\eqref{eq:prime-labels} satisfies \(S=-S\).
\end{enumerate}
\end{theorem}

\begin{proof}
\emph{(i).} If \(\iota A=A\), then \eqref{eq:reflection-pair} gives
\(P_A(\rho)=(-1)^p(P_A(-\rho)+s_pA(0))\). For even \(p\) the
polynomial \(A\) is odd, so \(A(0)=0\) and \(P_A\) is even. For odd
\(p\) the sum \(P_A(\rho)+P_A(-\rho)\) is the constant \(-s_pA(0)\).
In both cases \(P_A\) is reflection-symmetric.

Conversely, let \(P=P_A\) be reflection-symmetric. Suppose first that
\(p\) is even. Then \(w(z)=v_A(-z)\) solves \eqref{eq:main} with the
same operator, and it has a pole at \(i\pi\) because
\(-i\pi\equiv i\pi\) modulo \(2\pi i\). Its energy equals that of
\(v_A\), since every term of \eqref{eq:energyterms} and
\eqref{eq:energy} contains an even total number of derivatives.
By the fixed-energy Laurent uniqueness of
Section~\ref{sec:completeness}, \(w=v_A\). Letting \(z\to-\infty\) in
\(v_A(-z)=v_A(z)\) gives \(A(0)=0\), and \eqref{eq:reflection-pair}
becomes \(v_{\iota A}=v_A\). Hence \(\iota A=A\) by the uniqueness of
the form \eqref{eq:normalized-A}.
Suppose next that \(p\) is odd. By \eqref{eq:symmetric-class},
\(-P(-\rho)=P(\rho)-2P(0)\), and a direct substitution shows that
\(-v_A(-z)\) and \(v_A(z)+2P(0)\) both solve \eqref{eq:main} with the
operator \(P-2P(0)\). Both have a pole at \(i\pi\). For odd \(p\) the
Laurent expansion at a pole is unique, so
\(-v_A(-z)=v_A(z)+2P(0)\). By \eqref{eq:reflection-pair},
\(v_{\iota A}(z)=v_A(z)+2P(0)+s_pA(0)\). Letting \(z\to-\infty\) gives
\(s_pA(0)=-2P(0)\), and again \(v_{\iota A}=v_A\).

\emph{(ii).} On the subspace \(\{\alpha_j=0,\ j\text{ odd}\}\) we have
\(\iota A=A\), and \eqref{eq:reflection-S} shows that \(R_A\) is an odd
polynomial. The equations \(R_r=0\) with \(r\) even therefore hold
identically there. With \(m=\lfloor d/2\rfloor\), there remain the
\(m\) equations \(R_1,R_3,\ldots,R_{2m-1}\) in the \(m\) unknowns
\(\alpha_2,\alpha_4,\ldots,\alpha_{2m}\). By \eqref{eq:weights} their
weights are
\[
 \begin{array}{c|c}
 \text{variables}&2,4,\ldots,2m\\
 \text{equations}&2d,2d-2,\ldots,2(d-m+1).
 \end{array}
\]
Their highest weighted parts are the restrictions of those of the full
system. The highest weighted parts of the omitted equations vanish
identically on the subspace, because \(H_A\) is an odd polynomial when
\(\iota A=A\). A common zero of the restricted leading
forms is therefore a polynomial \(A\) with \(A\mid H_A\), and
\((\mathrm R_d)\) gives \(A=\rho^d\). The finite-cover argument of
Theorem~\ref{thm:length}, with \(\alpha_{2i}=y_i^{2i}\), gives
\[
 \len\Xp^{\mathrm{sym}}
 =\prod_{i=1}^{m}\frac{2(d-i+1)}{2i}=\binom dm .
\]

\emph{(iii).} A closed subscheme of a finite reduced scheme is reduced,
as in the proof of Proposition~\ref{cor:endpoints}, and a reduced point
of a square system is a simple solution of it. For even \(p\) the
polynomial \(A\) is odd, so \(A(0)=0\). For odd \(p\), a point with
\(A(0)=0\) would have \(P_A(0)=-s_pA(0)/2=0\) by the proof of~(i),
contrary to Proposition~\ref{cor:endpoints}. Finally, the map
\eqref{eq:iota} negates the roots of \(\overline A\), so it acts on the
labels by \(S\mapsto-S\), and the labelling is a bijection.
\end{proof}

The label count agrees with \eqref{eq:Fp}. The nonzero residues modulo
\(\ell=2d+1\) form \(d\) pairs \(\{s,-s\}\). A set \(S=-S\) of \(d\)
elements consists of \(d/2\) pairs if \(d\) is even, and of \(0\) and
\((d-1)/2\) pairs if \(d\) is odd, which gives
\(\binom{d}{\lfloor d/2\rfloor}\) sets in both cases. It contains
\(0\) exactly when \(d\) is odd, in agreement with the pulses and
fronts of part~(iii).

\begin{corollary}[Elliptic pairs with an even operator]
\label{cor:elliptic-symmetric}
Let \(p\) be even and assume \((\mathrm R_{p-1})\). On every fixed
nonsingular lattice, the elliptic pairs whose operator is an even
polynomial form a closed subscheme
\(\mathcal E_p^{\mathrm{sym}}(g_2,g_3)\) of
\(\mathcal E_p(g_2,g_3)\), defined by \(u_j=0\) for odd \(j\), and
\begin{equation}\label{eq:elliptic-symmetric-length}
 \len\mathcal E_p^{\mathrm{sym}}(g_2,g_3)=2F_p .
\end{equation}
If \(\Xp\) is reduced, \(\mathcal E_p^{\mathrm{sym}}(g_2,g_3)\) is
reduced for \((g_2,g_3)\) in a nonempty Zariski-open set of nonsingular
parameters.
\end{corollary}

\begin{proof}
If the profile \eqref{eq:elliptic-enumeration-profile} is even, the odd
part of \eqref{eq:main} is \(\sum_{j\text{ odd}}a_jD^jV=0\). The
functions \(D^jV\) have distinct pole orders, so \(a_j=0\) for odd
\(j\). Conversely, for an even operator \(V(-z)\) is a solution with
the same energy and the same pole, hence equals \(V\), as in the proof
of Theorem~\ref{thm:symmetric}.
On the subspace \(u_j=0\), \(j\) odd, the residual of Step~1 in the
proof of Theorem~\ref{thm:elliptic-count} is even. There remain the
equations and unknowns of even weight in
\eqref{eq:elliptic-weights}:
\[
 \begin{array}{c|c}
 \text{variables}&2,4,\ldots,p-2,p\\
 \text{equations}&p+2,p+4,\ldots,2p-2,2p.
 \end{array}
\]
The omitted leading forms vanish identically on the subspace, so the
remaining ones have only the trivial common zero, by Step~3 of that
proof. Weighted B\'ezout gives
\[
 \frac{(p+2)(p+4)\cdots(2p-2)\,(2p)}{2\cdot4\cdots(p-2)\,p}
 =2\binom{p-1}{p/2-1}=2F_p .
\]
At the node, reflection is \(Q\mapsto1-Q\), which fixes \(X\) and
changes the sign of \(Y\). An even nodal profile is
\(\beta+v_A\) with \(\iota A=A\), so the nodal fiber is
\eqref{eq:nodal-matching} with \(A\in\Xp^{\mathrm{sym}}\). If \(\Xp\)
is reduced, it has \(2F_p\) simple points by
Theorem~\ref{thm:symmetric} and Proposition~\ref{cor:endpoints}, and
Step~5 applies without change.
\end{proof}

For odd \(p\) there is no fixed-lattice analogue. The symmetric profile
is \(-P(0)\) plus a function odd about the pole, and the matching
system has one equation more than it has unknowns, so symmetric
elliptic pairs are not expected on a generic lattice. At \(p=3\) they
occur exactly when \(g_2=0\), as \eqref{eq:ks-elliptic-scheme} shows
for \(B=0\).

\begin{table}[htbp]
\centering
\caption{The symmetric counts: \(F_p\) rational-exponential pairs, and
\(2F_p\) elliptic pairs on a generic fixed lattice when \(p\) is even. They are lengths for \(p\le73\) and numbers of distinct simple
pairs under the reducedness hypothesis of
Theorem~\ref{thm:symmetric}.}
\label{tab:symmetric}
\begin{tabular}{@{}cllcc@{}}
\toprule
\(p\)&Symmetric operator&Example&\(F_p\)&Elliptic\\
\midrule
2&\(D^2+a_0\)&KdV&1&2\\
3&\(D^3+a_1D+a_0\)&KS without dispersion&2&--\\
4&\(D^4+a_2D^2+a_0\)&Kawahara&3&6\\
5&\(D^5+a_3D^3+a_1D+a_0\)&sixth-order model&6&--\\
6&\(D^6+a_4D^4+a_2D^2+a_0\)&seventh-order KdV&10&20\\
7&\(D^7+a_5D^5+a_3D^3+a_1D+a_0\)&&20&--\\
\bottomrule
\end{tabular}
\end{table}

Table~\ref{tab:symmetric} lists the first values. The rows \(p=2,3,4\)
are the KdV pulse of \eqref{eq:kdvb-pairs}, the two pairs with \(a=0\)
in Table~\ref{tab:ks}, and the Kawahara counts of
Section~\ref{sec:kawahara}. At \(p=5\) the scheme \(\mathcal X_5\) is
reduced (Section~\ref{sec:further}), so there are exactly six symmetric
pairs. With \(A=\rho^4+c_1\rho^2+c_2\) they are
\begin{equation}\label{eq:p5-symmetric}
 A=(\rho^2-1)(\rho^2-4),\qquad
 A=(\rho^2-1)(\rho^2-9),\qquad
 A=(\rho^2-1)\Bigl(\rho^2+\frac15\Bigr),
\end{equation}
and the three pairs with
\[
 1105c_1^3-2969c_1^2+5570c_1-4000=0,\qquad
 c_2=\frac{5c_1^3+21c_1-10}{79c_1+210}.
\]
The cubic has one real root. All six pairs are fronts, as
Theorem~\ref{thm:symmetric} requires, and the four with real
coefficients are, after the scales are restored, the solutions~(1.4)
of \citet{KudryashovMigita2007}. For instance,
\(A=(\rho^2-1)(\rho^2-4)\) gives
\[
 v_A=1260\bigl(40Q^3-60Q^4+24Q^5\bigr),\qquad
 P_A=\rho^5-55\rho^3+1254\rho-2520 .
\]

\subsection{Fixed equations}\label{sec:fixed}

The enumeration lets the operator vary. We now read it for a single
evolution equation
\begin{equation}\label{eq:pde-fixed}
 u_t+uu_x+\sum_{j=1}^{p}b_j\partial_x^{j+1}u=0,\qquad b_p\ne0,
\end{equation}
with fixed complex coefficients. Equation \eqref{eq:pde-fixed} is
invariant under translations and under the Galilean shifts
\(u\mapsto u+k\), \(x\mapsto x-kt\).

\begin{proposition}[The waves of a fixed equation]\label{prop:fixed}
Let \(\lambda\ne0\). Up to a translation and a Galilean shift, the
travelling waves \(u=w(x-ct)\) of \eqref{eq:pde-fixed} that are
rational in \(e^{\lambda(x-ct)}\), with one pole modulo
\(2\pi i/\lambda\), are
\begin{equation}\label{eq:fixed-wave}
 u=b_p\lambda^p\,v_A\bigl(\lambda(x-ct)\bigr),\qquad
 c=-b_p\lambda^pP_A(0),
\end{equation}
where \(A\) ranges over the points of \(\Xp\) that satisfy
\begin{equation}\label{eq:fixed-conditions}
 \frac{b_j}{b_p}=\lambda^{p-j}\,[\rho^j]P_A,\qquad 1\le j\le p-1 .
\end{equation}
The data \((A,\lambda)\) and \((\iota A,-\lambda)\) describe the same
wave.
\end{proposition}

\begin{proof}
Such a wave has finite limits as \(e^{\lambda(x-ct)}\to0\) and
\(\infty\), by Section~\ref{sec:systems}. A Galilean shift makes the
first limit zero, and the integration constant of
Section~\ref{sec:intro} is then zero. The substitution
\(w(\xi)=b_p\lambda^pv(\lambda\xi)\) turns the integrated equation into
\eqref{eq:main} with
\[
 [\rho^j]P=\frac{b_j}{b_p}\lambda^{j-p}\quad(1\le j\le p-1),\qquad
 P(0)=-\frac{c}{b_p\lambda^p}.
\]
After a translation, \(v\) has the normalization of this section, so
\(v=v_A\) and \(P=P_A\) by Proposition~\ref{prop:divisibility}.
The last assertion follows from \eqref{eq:reflection-pair}: the two
waves differ by the constant \(b_p\lambda^ps_pA(0)\), which is a
Galilean shift.
\end{proof}

Let \(\pi(A)=([\rho^{p-1}]P_A,\ldots,[\rho^{1}]P_A)\in\CC^{p-1}\), and
let \(\lambda\in\CC^\ast\) act on \(\CC^{p-1}\) with weights
\(1,2,\ldots,p-1\). Condition \eqref{eq:fixed-conditions} says that the
coefficient vector
\(\beta=(b_{p-1}/b_p,\ldots,b_1/b_p)\) lies in the orbit of \(\pi(A)\).
The rate rescaling moves an operator along a weighted curve, and this
is the reason for fixing the rate in the enumeration: \(N_p\) counts
these curves together with the waves that they carry.

\begin{corollary}[Solvable equations]\label{cor:solvable}
Assume that \(\Xp\) is finite, for instance that
\((\mathrm R_{p-1})\) holds.
\begin{enumerate}
\item[\textup{(i)}] Equation \eqref{eq:pde-fixed} has a travelling
wave that is rational in an exponential, with one pole per period, if
and only if \(\beta\) lies in one of the orbits
\(\CC^\ast\pi(A)\), \(A\in\Xp\). There are at most \((N_p+F_p)/2\)
such orbits.
\item[\textup{(ii)}] If \(\beta\ne0\), there are finitely many such
waves up to translations and Galilean shifts.
\item[\textup{(iii)}] If \(p\ge3\), the orbits form a set of dimension
at most one in \(\CC^{p-1}\). A generic equation
\eqref{eq:pde-fixed} has no such wave.
\end{enumerate}
\end{corollary}

\begin{proof}
Part~(i) restates Proposition~\ref{prop:fixed}. For the bound,
\(\pi(\iota A)=(-1)\cdot\pi(A)\) by \eqref{eq:reflection-pair}, so the
orbit depends only on the \(\iota\)-orbit of \(A\). The number of
\(\iota\)-orbits of points is half the number of points plus half the
number of fixed points, and these two numbers are at most
\(\len\Xp\le N_p\) and \(\len\Xp^{\mathrm{sym}}\le F_p\); the last
inequality is the B\'ezout bound in the proof of
Theorem~\ref{thm:symmetric}, which does not require
\((\mathrm R_{p-1})\). For~(ii), if \(\beta\ne0\) and
\(\beta\in\CC^\ast\pi(A)\), some coordinate \([\rho^j]P_A\) is nonzero,
and \eqref{eq:fixed-conditions} leaves at most \(p-j\) values of
\(\lambda\). Part~(iii) is a dimension count.
\end{proof}

When \(\beta=0\) the equation is invariant under the scaling itself,
and its waves come in one-parameter families of rates, as for the KdV
soliton. At \(p=2\) there are two orbits, \(\beta=0\) and
\(\beta\ne0\), and both carry waves: every KdV--Burgers equation has
the two fronts of Section~\ref{sec:kdvb}. At \(p=3\), with
\(\beta=(B,\mu)\), the orbits are the curves \(B^2=\sigma\mu\), and
Table~\ref{tab:ks} gives the four classical values
\begin{equation}\label{eq:ks-sigma}
 \sigma=\frac{16a^2}{a^2+19b}\in
 \Bigl\{0,\ 16,\ \frac{144}{47},\ \frac{256}{73}\Bigr\}.
\end{equation}
Here the bound \((N_3+F_3)/2=6\) is the number of \(\iota\)-orbits.
It is not attained, because the two orbits with \(a=0\) share
\(\sigma=0\), and \((\pm1,0)\) and \((\pm i,0)\) share \(\sigma=16\).
A floating-point evaluation of \(\pi\) on the exact solution lists
indicates \(17\) orbits at \(p=4\), where the three Kawahara pairs share
one orbit, and equality in the bound, \(66\) and \(236\) orbits, at
\(p=5\) and \(p=6\). We have not certified these three values.

For odd \(p\) the completeness theorem turns
Corollary~\ref{cor:solvable} into a statement about all meromorphic
travelling waves.

\begin{corollary}[Generic equations of odd order \(p\)]
\label{cor:generic}
Let \(p\ge5\) be odd and assume \((\mathrm R_{p-1})\). The coefficient
vectors \(\beta\in\CC^{p-1}\) for which \eqref{eq:pde-fixed} has a
nonconstant meromorphic travelling wave lie in an algebraic subset of
dimension at most two. In particular a generic equation has none.
For \(p=3\) this set is the union of the four curves
\(B^2=\sigma\mu\) of \eqref{eq:ks-sigma}.
\end{corollary}

\begin{proof}
By Theorem~\ref{thm:complete}, such a wave is rational, rational in an
exponential, or elliptic, with one pole modulo its periods. After a
Galilean shift that moves a constant solution to zero, the
substitution \(w=b_pv\) gives \eqref{eq:main} with
\([\rho^j]P=b_j/b_p\) for \(1\le j\le p-1\).
The rational-exponential waves give a set of dimension at most one, by
Corollary~\ref{cor:solvable}. For a rational wave
\eqref{eq:rational-system}, the argument of Step~3 in the proof of
Theorem~\ref{thm:elliptic-count} applies verbatim, with the simple-pole
term \(r_1z^{-1}\) included: it gives \(b=a_0=0\), and the Laplace
correspondence gives \(A\mid H_A\) for a monic \(A\) of degree
\(p-1\), now without the restriction \(A(0)=0\). By
\((\mathrm R_{p-1})\), \(A=\rho^{p-1}\), hence \(P=D^p\) and
\(\beta=0\).
For the elliptic waves, the family of all
\(\mathcal E_p(g_2,g_3)\) is finite over the \((g_2,g_3)\)-plane, by
Step~4 of the same proof. It has dimension two, and so has its image
under \(P\mapsto\beta\). For \(p=3\), Section~\ref{sec:ks} gives the
elliptic condition \(B^2=16\mu\), which is one of the four curves; the
point \(\beta=0\) lies on all of them.
\end{proof}

For even \(p\), Corollary~\ref{cor:solvable} and the dimension bound
for the elliptic pairs hold without change, but
Theorem~\ref{thm:complete} is available only when \(P\) is an even
polynomial.
